\documentclass[11pt,a4paper]{article}

\usepackage[utf8]{inputenc}
\usepackage[T1]{fontenc}
\usepackage[french]{babel}
\usepackage[a4paper,margin=2cm]{geometry}
\usepackage{amsmath,amssymb,siunitx}
\usepackage{lmodern}
\usepackage{enumitem}
\usepackage{fancyhdr}
\usepackage{lastpage}
\usepackage{xcolor}
\usepackage{amsthm}
\usepackage{array}
\usepackage{tabularx}

\sisetup{locale=FR}

\setlength{\headheight}{14pt}
\pagestyle{fancy}
\fancyhf{}
\fancyhead[L]{Terminale générale -- Spécialité Physique-Chimie}
\fancyhead[C]{Corrigé détaillé complet}
\fancyhead[R]{Page \thepage/\pageref{LastPage}}
\fancyfoot[C]{\textit{Corrigé détaillé avec rappels, commentaires physiques et réponses attendues}}

\newtheorem*{rappel}{Rappel}
\newtheorem*{remarque}{Remarque}
\newtheorem*{reponseattendue}{Réponse attendue}
\newtheorem*{commentaire}{Commentaire}

\begin{document}

\begin{center}
    {\Large \textbf{CORRIGÉ DÉTAILLÉ -- PHYSIQUE}}\\[0.3em]
    {\large \textbf{Sujet inspiré de GEIPI Polytech -- version approfondie corrigée}}\\[0.4em]
    {\large \textbf{Avec rappels de formules, commentaires physiques et réponses aux questions ouvertes}}
\end{center}

\vspace{0.5em}

\noindent\textbf{Corrections apportées au sujet :}
\begin{itemize}[leftmargin=1.8em]
    \item Dans l'exercice 1, on prend désormais l'interfrange
    \[
    i=\SI{2,5}{mm}
    \]
    afin d'obtenir une situation physiquement cohérente.
    \item Dans l'exercice 3, partie C, la vitesse
    \[
    V_0=\SI{22,0}{m.s^{-1}}
    \]
    est prise indépendamment des résultats précédents.
    \item Dans l'exercice 3, la portée est calculée dans un modèle simplifié où le skieur retombe à la même altitude que le point d'envol.
\end{itemize}

\vspace{0.5em}
\hrule
\vspace{0.6em}

\section*{Exercice 1 -- Optique ondulatoire : mesurer, interpréter, critiquer}

\noindent\textbf{Données :}
\[
\lambda=\SI{532}{nm},\qquad
D=\SI{2,50}{m},\qquad
c=\SI{3,00e8}{m.s^{-1}}.
\]

\vspace{0.3em}

\subsection*{Partie A -- Questions de base}

\begin{enumerate}[label=\textbf{1.\arabic*}, leftmargin=1.6em]

\item \textbf{Calcul de la fréquence du laser}

\begin{rappel}
Pour une onde lumineuse dans le vide :
\[
c=\lambda f
\qquad\Longrightarrow\qquad
f=\frac{c}{\lambda}.
\]
\end{rappel}

On convertit la longueur d’onde :
\[
\lambda=\SI{532}{nm}=532\times 10^{-9}\ \si{m}=5{,}32\times 10^{-7}\ \si{m}.
\]

Donc :
\[
f=\frac{3{,}00\times 10^8}{5{,}32\times 10^{-7}}
\simeq 5{,}64\times 10^{14}\ \si{Hz}.
\]

\[
\boxed{f \simeq 5{,}6\times 10^{14}\ \si{Hz}}
\]

\item \textbf{Nature de la lumière mise en évidence}

La diffraction met en évidence le \textbf{caractère ondulatoire} de la lumière.

Les interférences mettent elles aussi en évidence le \textbf{caractère ondulatoire} de la lumière.

\[
\boxed{\text{Dans les deux cas, on met en évidence la nature ondulatoire de la lumière.}}
\]

\item \textbf{Pourquoi un modèle purement corpusculaire est insuffisant}

Un modèle purement corpusculaire n’explique pas naturellement :
\begin{itemize}
    \item l’étalement d’un faisceau après passage dans une petite ouverture ;
    \item l’apparition de franges brillantes et sombres régulières.
\end{itemize}

Ces phénomènes sont caractéristiques d’ondes qui se propagent, se superposent et interfèrent.

\end{enumerate}

\subsection*{Partie B -- Mesure d’un diamètre de trou}

On observe une tache centrale de largeur :
\[
L=\SI{8,0}{mm}=8{,}0\times 10^{-3}\ \si{m}.
\]

\begin{enumerate}[label=\textbf{1.\arabic*}, resume, leftmargin=1.6em]

\item \textbf{Établissement de \(\theta \simeq \dfrac{L}{2D}\)}

\begin{rappel}
Pour les petits angles :
\[
\tan\theta \simeq \theta
\qquad (\theta \text{ en radians}).
\]
\end{rappel}

La demi-largeur de la tache vaut \(L/2\).  
Géométriquement :
\[
\tan\theta=\frac{L/2}{D}.
\]

Comme l’angle est petit :
\[
\theta \simeq \tan\theta \simeq \frac{L}{2D}.
\]

\[
\boxed{\theta \simeq \frac{L}{2D}}
\]

\item \textbf{Expression de \(a\)}

\begin{rappel}
Pour un trou circulaire :
\[
\theta \simeq \frac{1{,}22\,\lambda}{a}.
\]
\end{rappel}

En identifiant :
\[
\frac{L}{2D}\simeq \frac{1{,}22\,\lambda}{a}.
\]

Donc :
\[
a \simeq \frac{2{,}44\,D\lambda}{L}.
\]

\[
\boxed{a \simeq \frac{2{,}44\,D\lambda}{L}}
\]

\item \textbf{Calcul de \(a\)}

\[
a \simeq \frac{2{,}44\times 2{,}50\times 5{,}32\times 10^{-7}}{8{,}0\times 10^{-3}}.
\]

Numériquement :
\[
a \simeq 4{,}06\times 10^{-4}\ \si{m}.
\]

\[
\boxed{a \simeq 4{,}1\times 10^{-4}\ \si{m}}
\qquad\text{soit}\qquad
\boxed{a \simeq \SI{0,41}{mm}}
\]

\item \textbf{Commentaire physique}

On trouve un diamètre de l’ordre du dixième de millimètre.  
Cela correspond bien à une ouverture très petite à l’échelle usuelle, même si ce n’est pas microscopique au sens du micromètre.

\begin{commentaire}
Un très bon élève peut remarquer que \(\SI{0,41}{mm}\) est petit, mais pas \og microscopique \fg{} au sens strict. Cette remarque est pertinente et valorisable.
\end{commentaire}

\end{enumerate}

\subsection*{Partie C -- Mesure d’un espacement entre trous}

On corrige l’énoncé et on prend :
\[
i=\SI{2,5}{mm}=2{,}5\times 10^{-3}\ \si{m}.
\]

\begin{enumerate}[label=\textbf{1.\arabic*}, resume, leftmargin=1.6em]

\item \textbf{Calcul de \(b\)}

\begin{rappel}
Pour deux ouvertures distantes de \(b\) :
\[
i=\frac{\lambda D}{b}
\qquad\Longrightarrow\qquad
b=\frac{\lambda D}{i}.
\]
\end{rappel}

Donc :
\[
b=\frac{5{,}32\times 10^{-7}\times 2{,}50}{2{,}5\times 10^{-3}}
=5{,}32\times 10^{-4}\ \si{m}.
\]

\[
\boxed{b \simeq 5{,}3\times 10^{-4}\ \si{m}}
\qquad\text{soit}\qquad
\boxed{b \simeq \SI{0,53}{mm}}
\]

\item \textbf{Comparaison de \(a\) et \(b\)}

Nous avons :
\[
a \simeq \SI{0,41}{mm},
\qquad
b \simeq \SI{0,53}{mm}.
\]

Donc :
\[
\boxed{b>a}
\]

La situation est maintenant cohérente : les trous sont distincts et séparés par une distance supérieure à leur diamètre.

\item \textbf{Pourquoi les interférences sont modulées par une enveloppe de diffraction}

Chaque trou ne donne pas un simple rayon, mais une onde diffractée qui s’étale.

Les deux ondes issues des deux trous interfèrent entre elles : cela produit les franges.

Mais comme chacune des deux ondes est elle-même diffractée, l’intensité globale est limitée par une enveloppe plus large de diffraction.

\[
\boxed{\text{La diffraction détermine l'enveloppe, les interférences déterminent les franges fines.}}
\]

\end{enumerate}

\subsection*{Partie D -- Questions ouvertes}

\begin{enumerate}[label=\textbf{1.\arabic*}, resume, leftmargin=1.6em]

\item \textbf{Améliorer la précision de la mesure de \(a\)}

\begin{reponseattendue}
On attendait par exemple :
\begin{itemize}
    \item augmenter la distance \(D\) pour agrandir la tache ;
    \item utiliser un dispositif de mesure plus précis (caméra, capteur, analyse numérique) ;
    \item stabiliser le montage pour éviter les vibrations ;
    \item améliorer l’alignement du faisceau ;
    \item réduire la lumière parasite ;
    \item utiliser une source plus stable.
\end{itemize}
\end{reponseattendue}

\begin{commentaire}
Physiquement, une tache plus grande est souvent plus facile à mesurer, ce qui réduit l’incertitude relative sur \(L\).
\end{commentaire}

\item \textbf{Pourquoi ces modifications aident-elles ?}

\begin{reponseattendue}
Il fallait expliquer que :
\begin{itemize}
    \item une figure plus grande ou plus nette améliore la précision de lecture ;
    \item un montage stable évite le flou ;
    \item une meilleure source améliore le contraste ;
    \item un meilleur alignement évite les déformations géométriques.
\end{itemize}
\end{reponseattendue}

\item \textbf{Sens physique de la cohérence}

\begin{reponseattendue}
La cohérence signifie qu’il existe une relation de phase suffisamment stable entre les ondes pour former des interférences nettes et durables.  
Si cette relation varie trop rapidement, les franges se brouillent ou disparaissent.
\end{reponseattendue}

\item \textbf{Pourquoi un laser est mieux adapté qu’une lampe ?}

\begin{reponseattendue}
Parce qu’un laser émet une lumière :
\begin{itemize}
    \item très monochromatique ;
    \item très directive ;
    \item généralement très cohérente.
\end{itemize}
Cela permet d’obtenir des figures d’interférences bien visibles.
\end{reponseattendue}

\item \textbf{Causes de dégradation expérimentale}

\begin{reponseattendue}
Exemples attendus :
\begin{itemize}
    \item vibrations du montage ;
    \item turbulence de l’air ;
    \item mauvais alignement ;
    \item ouvertures mal usinées ;
    \item lumière parasite ;
    \item source insuffisamment monochromatique.
\end{itemize}
\end{reponseattendue}

\end{enumerate}

\vspace{0.6em}
\hrule
\vspace{0.6em}

\section*{Exercice 2 -- Refroidissement d’une boisson : du modèle simple à sa critique}

\noindent\textbf{Données :}
\[
\theta_{\mathrm{amb}}=\SI{20,0}{^\circ C},
\qquad
\theta(0)=\SI{85,0}{^\circ C},
\qquad
k=2{,}1\times 10^{-4}\ \si{s^{-1}}.
\]

Le modèle est :
\[
\frac{d\theta}{dt}=-k\big(\theta(t)-\theta_{\mathrm{amb}}\big).
\]

\subsection*{Partie A -- Lecture physique du modèle}

\begin{enumerate}[label=\textbf{2.\arabic*}, leftmargin=1.6em]

\item \textbf{Analyse d’unités}

Le membre de gauche \(\dfrac{d\theta}{dt}\) s’exprime en \(\si{^\circ C.s^{-1}}\) ou \(\si{K.s^{-1}}\).

Le terme \(\theta-\theta_{\mathrm{amb}}\) s’exprime en \(\si{^\circ C}\) ou \(\si{K}\).

Donc :
\[
[k]=\si{s^{-1}}.
\]

\[
\boxed{[k]=\si{s^{-1}}}
\]

\item \textbf{Signe de la dérivée au début}

Au départ :
\[
\theta(0)-\theta_{\mathrm{amb}}=85-20=65>0.
\]

Donc :
\[
\frac{d\theta}{dt}<0.
\]

La température décroît : la boisson refroidit.

\item \textbf{À l’équilibre}

Si
\[
\theta(t)=\theta_{\mathrm{amb}},
\]
alors :
\[
\frac{d\theta}{dt}=0.
\]

\[
\boxed{\text{À l’équilibre thermique, la température devient constante.}}
\]

\item \textbf{Pourquoi ce modèle est raisonnable}

Plus l’écart de température avec l’air ambiant est grand, plus les échanges thermiques sont intenses.  
Le modèle suppose que la vitesse de refroidissement est proportionnelle à cet écart, ce qui est une hypothèse très classique et raisonnable.

\end{enumerate}

\subsection*{Partie B -- Résolution}

\begin{rappel}
La solution admise est :
\[
\theta(t)=\theta_{\mathrm{amb}}+\big(\theta(0)-\theta_{\mathrm{amb}}\big)e^{-kt}.
\]
\end{rappel}

Ici :
\[
\theta(t)=20+65\,e^{-kt}.
\]

\begin{enumerate}[label=\textbf{2.\arabic*}, resume, leftmargin=1.6em]

\item \textbf{Température au bout de 10 min}

\[
10\ \text{min} = 600\ \text{s}.
\]

Donc :
\[
\theta(600)=20+65e^{-2{,}1\times 10^{-4}\times 600}
=20+65e^{-0{,}126}.
\]

Or :
\[
e^{-0{,}126}\simeq 0{,}882.
\]

Donc :
\[
\theta(600)\simeq 20+65\times 0{,}882 \simeq 77{,}3.
\]

\[
\boxed{\theta(10\ \text{min})\simeq \SI{77,3}{^\circ C}}
\]

\item \textbf{Temps pour atteindre \(\SI{55}{^\circ C}\)}

On résout :
\[
55=20+65e^{-kt}.
\]

Donc :
\[
35=65e^{-kt}
\qquad\Longrightarrow\qquad
e^{-kt}=\frac{7}{13}.
\]

En prenant le logarithme :
\[
-kt=\ln\left(\frac{7}{13}\right).
\]

Ainsi :
\[
t=-\frac{1}{k}\ln\left(\frac{7}{13}\right).
\]

Numériquement :
\[
t\simeq -\frac{1}{2{,}1\times 10^{-4}}\times \ln\left(\frac{7}{13}\right)
\simeq 2{,}95\times 10^3\ \si{s}.
\]

\[
\boxed{t\simeq \SI{2950}{s}\simeq \SI{49}{min}}
\]

\item \textbf{Limite quand \(t\to +\infty\)}

Comme
\[
e^{-kt}\to 0,
\]
on obtient :
\[
\theta(t)\to \theta_{\mathrm{amb}}.
\]

\[
\boxed{\lim_{t\to +\infty}\theta(t)=\SI{20,0}{^\circ C}}
\]

Interprétation : à long terme, la boisson prend la température ambiante.

\end{enumerate}

\subsection*{Partie C -- Constante de temps}

On pose :
\[
\tau=\frac{1}{k}.
\]

\begin{enumerate}[label=\textbf{2.\arabic*}, resume, leftmargin=1.6em]

\item \textbf{Calcul de \(\tau\)}

\[
\tau=\frac{1}{2{,}1\times 10^{-4}}
\simeq 4{,}76\times 10^3\ \si{s}.
\]

\[
\boxed{\tau\simeq \SI{4,8e3}{s}}
\qquad\text{soit environ}\qquad
\boxed{\tau\simeq \SI{79}{min}}
\]

\item \textbf{Valeur de \(\theta(\tau)\)}

\[
\theta(\tau)=\theta_{\mathrm{amb}}+\big(\theta(0)-\theta_{\mathrm{amb}}\big)e^{-k\tau}.
\]

Or :
\[
k\tau=1.
\]

Donc :
\[
\theta(\tau)=\theta_{\mathrm{amb}}+\big(\theta(0)-\theta_{\mathrm{amb}}\big)e^{-1}.
\]

\[
\boxed{
\theta(\tau)=\theta_{\mathrm{amb}}+\frac{\theta(0)-\theta_{\mathrm{amb}}}{e}
}
\]

\item \textbf{Interprétation physique}

Au temps \(\tau\), l’écart à la température ambiante a été divisé par \(e\), c’est-à-dire environ par \(2{,}72\).

\item \textbf{Effet d’une meilleure isolation}

Une tasse mieux isolée ralentit les échanges thermiques.

Donc le refroidissement est plus lent, ce qui correspond à un \(k\) plus petit, donc à une constante de temps plus grande.

\[
\boxed{\text{Meilleure isolation} \Longrightarrow \tau \text{ plus grande}}
\]

\end{enumerate}

\subsection*{Partie D -- Questions ouvertes}

\begin{enumerate}[label=\textbf{2.\arabic*}, resume, leftmargin=1.6em]

\item \textbf{Pourquoi l’évaporation accélère le refroidissement ?}

\begin{reponseattendue}
L’évaporation emporte de l’énergie thermique : les molécules qui quittent le liquide emportent de l’énergie.  
Cela provoque une perte d’énergie supplémentaire et donc un refroidissement plus rapide.
\end{reponseattendue}

\item \textbf{Le coefficient \(k\) est-il vraiment constant ?}

\begin{reponseattendue}
Pas forcément. Le coefficient \(k\) modélise globalement les échanges thermiques, mais ceux-ci peuvent dépendre :
\begin{itemize}
    \item de la température ;
    \item des mouvements d’air ;
    \item de l’évaporation ;
    \item de la géométrie de la tasse ;
    \item de la surface libre du liquide.
\end{itemize}
Donc \(k\) constant est une approximation utile, mais simplificatrice.
\end{reponseattendue}

\item \textbf{« Si on met deux fois plus de boisson, le temps de refroidissement double exactement »}

\begin{reponseattendue}
Il fallait répondre que :
\begin{itemize}
    \item qualitativement, plus de liquide signifie plus d’énergie thermique à perdre, donc un refroidissement plus lent ;
    \item mais quantitativement, le temps ne double pas forcément exactement.
\end{itemize}

En effet, si la masse de liquide change, l’inertie thermique change, mais la surface d’échange avec l’air ne double pas nécessairement. Les mécanismes de transfert thermique ne suivent donc pas une simple proportionnalité exacte.
\end{reponseattendue}

\[
\boxed{\text{Affirmation plausible qualitativement, mais fausse si on dit “exactement”.}}
\]

\end{enumerate}

\vspace{0.6em}
\hrule
\vspace{0.6em}

\section*{Exercice 3 -- Mouvement d’un skieur : énergie, trajectoire, frottements}

\noindent\textbf{Données :}
\[
m=\SI{60}{kg},
\qquad
g=\SI{9,81}{m.s^{-2}}.
\]

\subsection*{Partie A -- Descente sans frottements}

On part du repos depuis une hauteur :
\[
h=\SI{35}{m}.
\]

\begin{enumerate}[label=\textbf{3.\arabic*}, leftmargin=1.6em]

\item \textbf{Expression de \(V_E\)}

\begin{rappel}
Sans frottements, l’énergie mécanique se conserve :
\[
E_{m,\text{initiale}}=E_{m,\text{finale}}.
\]
\end{rappel}

Au départ :
\[
E_{c,D}=0,
\qquad
E_{p,D}=mgh.
\]

Au point \(E\), on prend l’énergie potentielle nulle :
\[
E_{c,E}=\frac12 mV_E^2.
\]

Donc :
\[
mgh=\frac12 mV_E^2.
\]

Ainsi :
\[
V_E=\sqrt{2gh}.
\]

\[
\boxed{V_E=\sqrt{2gh}}
\]

\item \textbf{Calcul de la vitesse}

\[
V_E=\sqrt{2\times 9{,}81\times 35}
=\sqrt{686{,}7}
\simeq 26{,}2\ \si{m.s^{-1}}.
\]

\[
\boxed{V_E\simeq \SI{26,2}{m.s^{-1}}}
\]

En km/h :
\[
26{,}2\times 3{,}6\simeq 94{,}3\ \si{km.h^{-1}}.
\]

\[
\boxed{V_E\simeq \SI{94}{km.h^{-1}}}
\]

\item \textbf{Énergie cinétique au point \(E\)}

\[
E_c=\frac12 mV_E^2.
\]

Donc :
\[
E_c\simeq \frac12\times 60\times 686{,}7
\simeq 2{,}06\times 10^4\ \si{J}.
\]

\[
\boxed{E_c\simeq 2{,}1\times 10^4\ \si{J}}
\]

\item \textbf{Commentaire}

La vitesse est très élevée, ce qui est cohérent avec un saut à ski.  
L’énergie cinétique est importante : cela montre qu’un faible changement de vitesse ou d’orientation peut avoir un effet notable.

\end{enumerate}

\subsection*{Partie B -- Avec dissipation d’énergie}

On suppose qu’une énergie mécanique
\[
E_f=\SI{3,5e3}{J}
\]
est dissipée.

\begin{enumerate}[label=\textbf{3.\arabic*}, resume, leftmargin=1.6em]

\item \textbf{Bilan énergétique}

\[
mgh=\frac12 mV_E^2+E_f.
\]

\[
\boxed{mgh=\frac12 mV_E^2+E_f}
\]

\item \textbf{Nouvelle vitesse}

\[
\frac12 mV_E^2=mgh-E_f.
\]

D’abord :
\[
mgh=60\times 9{,}81\times 35 \simeq 2{,}06\times 10^4\ \si{J}.
\]

Donc :
\[
mgh-E_f \simeq 2{,}06\times 10^4 - 3{,}5\times 10^3
=1{,}71\times 10^4\ \si{J}.
\]

Puis :
\[
V_E^2=\frac{2\times 1{,}71\times 10^4}{60}\simeq 570.
\]

Ainsi :
\[
V_E\simeq \sqrt{570}\simeq 23{,}9\ \si{m.s^{-1}}.
\]

\[
\boxed{V_E\simeq \SI{23,9}{m.s^{-1}}}
\]

\item \textbf{Comparaison}

Sans pertes :
\[
V_E\simeq \SI{26,2}{m.s^{-1}}.
\]

Avec pertes :
\[
V_E\simeq \SI{23,9}{m.s^{-1}}.
\]

L’écart est d’environ :
\[
2{,}3\ \si{m.s^{-1}}.
\]

\item \textbf{Conséquences possibles}

Oui, cet écart peut être très important car la portée, la hauteur maximale et la sécurité du saut dépendent fortement de la vitesse d’envol.

\end{enumerate}

\subsection*{Partie C -- Vol dans l’air}

\noindent\textbf{Dans cette partie, on prend indépendamment des résultats précédents :}
\[
V_0=\SI{22,0}{m.s^{-1}},
\qquad
\alpha=11^\circ.
\]

On néglige l’action de l’air.

\begin{enumerate}[label=\textbf{3.\arabic*}, resume, leftmargin=1.6em]

\item \textbf{Équations horaires}

\begin{rappel}
Sans frottements :
\begin{itemize}
    \item sur l’axe horizontal : mouvement uniforme ;
    \item sur l’axe vertical : mouvement uniformément accéléré d’accélération \(-g\).
\end{itemize}
\end{rappel}

Les composantes de la vitesse initiale sont :
\[
V_{0x}=V_0\cos\alpha,
\qquad
V_{0z}=V_0\sin\alpha.
\]

Donc :
\[
\boxed{x(t)=V_0\cos\alpha\, t}
\]
et
\[
\boxed{z(t)=V_0\sin\alpha\, t-\frac12 gt^2}
\]

\item \textbf{Équation de la trajectoire}

On a :
\[
t=\frac{x}{V_0\cos\alpha}.
\]

En remplaçant dans \(z(t)\) :
\[
z=x\tan\alpha-\frac{g}{2V_0^2\cos^2\alpha}x^2.
\]

\[
\boxed{
z(x)=x\tan\alpha-\frac{g}{2V_0^2\cos^2\alpha}x^2
}
\]

C’est donc une parabole.

\item \textbf{Hauteur maximale}

La vitesse verticale vaut :
\[
v_z(t)=V_0\sin\alpha-gt.
\]

Au sommet :
\[
v_z=0
\qquad\Longrightarrow\qquad
t_s=\frac{V_0\sin\alpha}{g}.
\]

La hauteur maximale vaut :
\[
z_{\max}=\frac{(V_0\sin\alpha)^2}{2g}.
\]

Or :
\[
\sin 11^\circ \simeq 0{,}191.
\]

Donc :
\[
V_0\sin\alpha \simeq 22\times 0{,}191 \simeq 4{,}20\ \si{m.s^{-1}}.
\]

Ainsi :
\[
z_{\max}\simeq \frac{4{,}20^2}{2\times 9{,}81}
=\frac{17{,}6}{19{,}62}
\simeq 0{,}90\ \si{m}.
\]

\[
\boxed{z_{\max}\simeq \SI{0,90}{m}}
\]

\item \textbf{Durée du vol dans le modèle simplifié}

On suppose ici que le skieur retombe à la même altitude que le point d’envol.

Il faut résoudre :
\[
z(t)=0
\]
avec \(t\neq 0\), soit :
\[
V_0\sin\alpha\, t-\frac12 gt^2=0.
\]

Donc :
\[
t_f=\frac{2V_0\sin\alpha}{g}.
\]

Numériquement :
\[
t_f=\frac{2\times 4{,}20}{9{,}81}\simeq 0{,}86\ \si{s}.
\]

\[
\boxed{t_f\simeq \SI{0,86}{s}}
\]

\item \textbf{Portée horizontale}

\[
x_f=V_0\cos\alpha\, t_f.
\]

Or :
\[
\cos 11^\circ \simeq 0{,}982,
\qquad
V_0\cos\alpha \simeq 22\times 0{,}982 \simeq 21{,}6\ \si{m.s^{-1}}.
\]

Donc :
\[
x_f\simeq 21{,}6\times 0{,}86 \simeq 18{,}6\ \si{m}.
\]

\[
\boxed{x_f\simeq \SI{18,6}{m}}
\]

\end{enumerate}

\subsection*{Partie D -- Frottement fluide simplifié}

On admet un modèle très simplifié avec une force de frottement :
\[
f=kv
\]
opposée au mouvement.

\begin{enumerate}[label=\textbf{3.\arabic*}, resume, leftmargin=1.6em]

\item \textbf{Pourquoi cette force diminue-t-elle la portée ?}

\begin{reponseattendue}
La force de frottement est opposée au mouvement : elle retire de l’énergie mécanique au système et fait diminuer la vitesse.  
Comme la composante horizontale de la vitesse diminue elle aussi, la portée est plus faible.
\end{reponseattendue}

\item \textbf{Agit-elle sur la norme seulement ou aussi sur la direction ?}

\begin{reponseattendue}
Elle agit sur les deux.  
Comme elle est opposée au vecteur vitesse, elle modifie sa norme, mais aussi la direction du mouvement au cours du temps.
\end{reponseattendue}

\item \textbf{Effet plus marqué au début ou à la fin ?}

\begin{reponseattendue}
Si \(f=kv\), alors plus la vitesse est grande, plus le frottement est fort.  
L’effet de l’air est donc généralement plus marqué lorsque la vitesse est plus grande, souvent au début du vol.
\end{reponseattendue}

\item \textbf{Position aérodynamique}

\begin{reponseattendue}
Une posture plus aérodynamique réduit la résistance de l’air.  
Le skieur conserve donc mieux sa vitesse, perd moins d’énergie et peut augmenter sa portée.
\end{reponseattendue}

\end{enumerate}

\subsection*{Partie E -- Question ouverte de synthèse}

\begin{enumerate}[label=\textbf{3.\arabic*}, resume, leftmargin=1.6em]

\item \textbf{Influence des paramètres}

\begin{reponseattendue}
On attendait une analyse structurée :
\begin{itemize}
    \item \textbf{hauteur de départ :} elle augmente l’énergie potentielle initiale, donc potentiellement la vitesse d’envol ;
    \item \textbf{frottements sur le tremplin :} s’ils diminuent, moins d’énergie est dissipée ;
    \item \textbf{angle d’envol :} il règle le compromis entre hauteur gagnée et portée ;
    \item \textbf{aérodynamique :} elle conditionne les pertes en vol.
\end{itemize}
\end{reponseattendue}

\item \textbf{Pourquoi il n’existe pas un “paramètre miracle”}

\begin{reponseattendue}
Parce que les paramètres agissent ensemble et parfois en sens opposés.  
Par exemple :
\begin{itemize}
    \item un angle trop grand peut augmenter la hauteur mais réduire la portée ;
    \item une grande vitesse sans bonne aérodynamique peut être mal exploitée ;
    \item une bonne posture ne compense pas totalement une vitesse initiale insuffisante.
\end{itemize}
La performance résulte donc d’un compromis physique global.
\end{reponseattendue}

\[
\boxed{\text{La performance dépend d’un équilibre entre plusieurs paramètres, pas d’un seul.}}
\]

\end{enumerate}

\vspace{0.8em}
\begin{center}
\textbf{Fin du corrigé}
\end{center}

\end{document}